% !TEX root = ../thesis.tex

%=========================================================
\section{Dynamical Coulomb Blockade}
\label{sec:ssn-set:dynamic}
%=========================================================

    SET2 provided a contrasting outcome to the gate controlled transport of SET1. I first characterized the device before opening the mechanically controlled break junction. I then partially opened the junction while retaining a contact configuration with a conductance of several $G_0$. Gate sweeps of the intact superconducting device and of the partially opened device in both the normal and superconducting states showed no resolvable modulation. These measurements covered the same applied gate voltage range in which SET1 showed periodic transport. I therefore found no evidence that opening SET2 produced a gate controlled island. Instead, the partially opened configuration developed a gate independent zero bias anomaly consistent with dynamical Coulomb blockade.

    %=========================================================
    \subsection*{SIN Tunnel Barrier}
    %=========================================================

    \singlefitfigure{ssn-set/ssn2-sin}{
        Superconducting transport of SET2 before opening the MCBJ. The measured \textit{I--V}, \textit{dI--dV}, and \textit{dV--dI} responses are compared with an SIN BCS Dynes fit. The fit reproduces the quasiparticle onset and establishes the normal-state resistance and superconducting parameters of the fixed tunnel barrier.
    }{fig:ssn2-set:sin}

    Before opening the break junction, its resistance was much smaller than that of the fixed Al--AlO\textsubscript{x}--Cu tunnel barrier. The device therefore behaved approximately as a single SIN tunnel junction. Figure~\ref{fig:ssn2-set:sin} shows its superconducting transport measured at $T_\mathrm{bath}=\qty{36(1)}{\milli\kelvin}$ together with a fit using the BCS density of states with Dynes broadening. The fit yielded the normal-state resistance and conductance
    \begin{align}
        R_\mathrm{TB} = \qty{32.379(2)}{\kilo\ohm}\,, &&
        G_\mathrm{TB} = 0.39859(3)\,G_0\,,
        \label{eq:ssn-set:dynamic:intact-normal}
    \end{align}
    and the superconducting parameters
    \begin{align}
        \Delta_0 = \qty{195.99(2)}{\micro\electronvolt}\,, &&
        T_\mathrm{fit} = \qty{59.5(8)}{\milli\kelvin}\,, &&
        \gamma = \qty{1.48(4)}{\micro\electronvolt}\,.
        \label{eq:ssn-set:dynamic:intact-superconducting}
    \end{align}
    After superconductivity was suppressed, the intact device showed no resolvable deviation from a linear normal-state response and no zero bias anomaly within the measurement noise. This measurement established the fixed tunnel barrier as a reference before the mechanical configuration was changed.

    %=========================================================
    \subsection*{Normal-State Dynamical Coulomb Blockade}
    %=========================================================

    I then partially opened the MCBJ. Unlike SET1, the resulting configuration did not show periodic gate modulation. Instead, its normal-state conductance measured at $T_\mathrm{bath}=\qty{70(1)}{\milli\kelvin}$ showed a gate independent suppression around zero bias, as shown in Fig.~\ref{fig:ssn2-set:dcb}. This behavior motivated a DCB fit. I fitted the complete measured \textit{I--V} range using the DCB formulation of \textit{P(E)} Theory from Sec.~\ref{subsec:stochastic:dcb} and the same effective parallel-RC environment ansatz used in Sec.~\ref{subsec:ac:photon-assisted-transport}. The normal-state resistance, effective resistance, effective capacitance, and effective temperature were free parameters in a single simultaneous fit. The fit reproduced the voltage dependent conductance with
    \begin{align}
        R_\mathrm{N} = \qty{34.43(3)}{\kilo\ohm}\,, &&
        G_\mathrm{N} = 0.37488(3)\,G_0\,,
        \label{eq:ssn-set:dynamic:dcb-normal}
    \end{align}
    and the effective environmental parameters
    \begin{align}
        R_\mathrm{eff} = \qty{0.446(3)}{\kilo\ohm}\,, &&
        C_\mathrm{eff} = \qty{4.36(5)}{\femto\farad}\,, &&
        T_\mathrm{eff} = \qty{77(6)}{\milli\kelvin}\,.
        \label{eq:ssn-set:dynamic:dcb-environment}
    \end{align}
    These parameters define the RC cutoff energy
    \begin{align}
        E_\mathrm{RC} = \frac{\hbar}{R_\mathrm{eff}C_\mathrm{eff}} = \qty{339(5)}{\micro\electronvolt}\,.
        \label{eq:ssn-set:dynamic:dcb-cutoff}
    \end{align}
    The contrast with the linear response of the intact device linked the anomaly to the mechanically modified configuration. Together with the absence of gate modulation, the effective RC fit supported an environmental DCB interpretation rather than static Coulomb blockade of a gate controlled island. I treated $R_\mathrm{eff}$, $C_\mathrm{eff}$, and the resulting $E_\mathrm{RC}$ as phenomenological parameters because the measurement did not identify them with individual circuit elements. 

    The measured bath temperature remained stable during the measurement, and I resolved no corresponding cryostat temperature drift. No temperature series or independent characterization of the effective environment was available. The RC fit therefore did not uniquely distinguish DCB from other contact or density of states effects, and stable bath temperature did not exclude electronic heating. I consequently used DCB as an effective phenomenological interpretation of the gate independent anomaly rather than as a uniquely established microscopic mechanism.
    
    Subtracting the intact tunnel barrier resistance from the normal-state resistance of the partially opened configuration gave an estimate for the normal-state break junction resistance,
    \begin{align}
        R_\mathrm{BJ} = \qty{2.05(3)}{\kilo\ohm}\,, &&
        G_\mathrm{BJ} = 6.30(9)\,G_0\,.
        \label{eq:ssn-set:dynamic:break-junction-normal}
    \end{align}
    This estimate showed that the MCBJ remained a high-conductance multichannel contact rather than entering the tunnel regime. Its conductance was compatible with an atomic-scale metallic constriction, although the transport measurement did not determine the contact geometry or number of atoms.

    \singlefitfigure{ssn-set/ssn2-dcb}{
        Normal-state transport of SET2 after partially opening the MCBJ. The measured \textit{I--V}, \textit{dI--dV}, and \textit{dV--dI} responses are compared with the dynamical Coulomb blockade fit for an effective RC environment. The gate independent zero bias conductance suppression contrasts with the linear normal-state response of the intact device.
    }{fig:ssn2-set:dcb}

    %=========================================================
    \subsection*{Superconducting Crossover}
    %=========================================================

    Directly after the normal-state measurement, I restored superconductivity and measured the device again at $T_\mathrm{bath}=\qty{81(1)}{\milli\kelvin}$ without mechanically adjusting the MCBJ. The two measurements therefore probed the same nominal contact configuration back to back. In the superconducting state, however, a single SIN BCS Dynes model did not reproduce the complete voltage range, as shown in Fig.~\ref{fig:ssn2-set:sin-dcb}. I consequently analyzed a low bias region (LB), $\lvert V\rvert<\qty{0.50}{\milli\volt}$, and a high bias region (HB), $\lvert V\rvert>\qty{0.55}{\milli\volt}$, separately.

    \singlefitfigure{ssn-set/ssn2-sin-dcb}{
        Superconducting transport of the partially opened SET2 configuration measured directly after the normal-state response in Fig.~\ref{fig:ssn2-set:dcb}. The panels compare the measured \textit{I--V}, \textit{dI--dV}, and \textit{dV--dI} responses with the LB Dynes fit and the HB reference calculated with the same superconducting parameters. The derivative responses resolve the crossover from the LB conductance scale toward the HB scale.
    }{fig:ssn2-set:sin-dcb}

    In the LB region, the response resembled that of an SIN tunnel junction. A Dynes fit yielded
    \begin{align}
        R_\mathrm{N}^\mathrm{LB} = \qty{38.34(2)}{\kilo\ohm}\,, &&
        G_\mathrm{N}^\mathrm{LB} = 0.33657(2)\,G_0\,,
        \label{eq:ssn-set:dynamic:lb-normal}
    \end{align}
    and the superconducting parameters
    \begin{align}
        \Delta_0 = \qty{195.94(6)}{\micro\electronvolt}\,, &&
        T_\mathrm{fit} = \qty{121(2)}{\milli\kelvin}\,, &&
        \gamma = \qty{0.88(9)}{\micro\electronvolt}\,.
        \label{eq:ssn-set:dynamic:lb-superconducting}
    \end{align}
    In the HB region, the measured differential conductance approached a larger value. I fitted the \textit{dI--dV} response in this region while keeping $\Delta_0$, $T_\mathrm{fit}$, and $\gamma$ fixed to the values extracted from the LB fit. The remaining conductance scale was
    \begin{align}
        R_\mathrm{N}^\mathrm{HB} = \qty{32.13(2)}{\kilo\ohm}\,, &&
        G_\mathrm{N}^\mathrm{HB} = 0.40172(3)\,G_0\,.
        \label{eq:ssn-set:dynamic:hb-normal}
    \end{align}
    The HB conductance was approximately $19\%$ larger than the conductance inferred from the LB response. Its resistance differed from the intact tunnel barrier resistance by only
    \begin{align}
        \left|R_\mathrm{N}^\mathrm{HB}-R_\mathrm{TB}\right| = \qty{0.25(2)}{\kilo\ohm}\,,
        \label{eq:ssn-set:dynamic:hb-barrier-difference}
    \end{align}
    corresponding to a relative difference of $< 1\%$. Thus, the HB resistance closely matched the independently measured tunnel barrier resistance.

    The derivative panels resolved the crossover between the LB and HB responses at
    \begin{align}
        V_\mathrm{x} = \qty{0.52(3)}{\milli\volt}\,, &&
        I_\mathrm{x} = \qty{13(1)}{\nano\ampere}\,.
        \label{eq:ssn-set:dynamic:crossover}
    \end{align}
    I identified the crossover interval from the visible changes in \textit{dI--dV} and \textit{dV--dI} without using the RC energy scale. I approximated the uncertainties by the half-widths of these intervals. They therefore quantify the crossover width rather than statistical fit uncertainty.

    The crossover provided an additional comparison with the normal-state DCB fit. Its bias energy and the sum of the superconducting gap and RC cutoff energy were
    \begin{align}
        eV_\mathrm{x} = \qty{520(30)}{\micro\electronvolt}\,, &&
        \Delta_0+E_\mathrm{RC} = \qty{535(5)}{\micro\electronvolt}\,.
        \label{eq:ssn-set:dynamic:crossover-energy}
    \end{align}
    The two scales agreed within the estimated crossover width. This correspondence was consistent with a crossover from an environmentally modified quasiparticle response toward the bare SIN response once the excess bias energy became comparable to the RC cutoff energy. Because the effective RC environment produced a gradual energy dependence rather than a sharp threshold, I did not identify this agreement as a derived onset condition.

    Expressed as resistance, the separation between the two fitted scales was
    \begin{align}
        \Delta R_\mathrm{N} = R_\mathrm{N}^\mathrm{LB}-R_\mathrm{N}^\mathrm{HB} = \qty{6.21(3)}{\kilo\ohm}\,.
        \label{eq:ssn-set:dynamic:resistance-difference}
    \end{align}
    This value was approximately \qty{4}{\percent} below $h/(4e^2)\approx\qty{6.45}{\kilo\ohm}$. This scale can be written as the resistance quantum for charge-$2e$ transfer, but it also equals the resistance of two parallel spin-degenerate unit-transmission electron channels. The present treatment within \textit{P(E)} Theory did not predict an additive resistance of this magnitude, so I retained the proximity as a numerical observation without assigning either interpretation.

    SET2 thus showed a gate independent zero-bias anomaly rather than the static gate controlled blockade observed for SET1. The normal-state response was consistent with DCB in an effective RC environment. In the superconducting state, the HB resistance agreed with the independently measured tunnel barrier resistance, while the LB Dynes fit yielded a larger effective resistance. The crossover energy agreed with $\Delta_0+E_\mathrm{RC}$ within its estimated width, while the difference between the fitted resistance scales was close to $R_\mathrm{Q}^{2e}$. The first correspondence was consistent with environmental modification of the quasiparticle response. The second did not follow directly from the model used here. I therefore report both without assigning a microscopic mechanism to the LB--HB crossover.

    Taken together, the two devices showed that the same nominal SSN design accessed distinct transport regimes when the MCBJ was opened to different contact conductances. SET1 sampled the weak to intermediate coupling regime. Its normal-state response was described at a minimal level by orthodox SET theory, and it showed gate dependent switching in the superconducting state. SET2 remained in the strong coupling regime and developed a gate independent zero-bias anomaly consistent with an effective DCB description. This contrast is consistent with the broader coupling regime behavior reported for the same SSN architecture by Laura Sobral Rey \textit{et al.}~\cite{sobral_rey_hybrid_sset_2024,sobral_rey_interplay_2022}. The measurements did not identify the microscopic contact properties responsible for these different outcomes. Because each device provided only one usable configuration, the comparison established the observed regimes but not a systematic transition between them.
